\subsection{20.122(a)--(c): minimum intersections outside the Fitting subgroup}
\label{cand:20.122}
\problemstatement{20.122}
For nilpotent \(A,B,C\le G\), the question defines
\(\operatorname{Min}_G(A,B,C)\) using all inclusion-minimal
intersections \(A\cap B^x\cap C^y\), and \(\min_G(A,B,C)\) using
those of minimum order. Both can lie outside \(F(G)\), even for soluble
\(G\). The distinction matters: the experiment's earlier smaller
example addressed only inclusion-minimal intersections.
Zenkov's earlier examples separate the analogous notions for pairs
of subgroups \cite{zenkov-intersections}. The argument below
computes the full triple-intersection family for a different group.

Let \(\Omega=\F_3^2\), \(o=0\), and let \(H\cong D_8\) be the eight
signed coordinate permutations. It is generated by
\(d(x,y)=(-x,y)\) and \(s(x,y)=(y,x)\).
Put
\[
 Q=\F_3^2\rtimes H,\qquad
 E=C_2^\Omega=\langle z_w:w\in\Omega\rangle,\qquad
 G=E\rtimes Q.
\]
In its natural action on \(\Omega\times\F_2\), a base element flips
selected fibers and \(Q\) permutes their indices.
This soluble group has order \(512\cdot72=36864\).
Take
\[
 Z=\langle z_o\rangle,\qquad
 A=Z\times H,\qquad B=C=G_{(o,0)}.
\]
The orders of \(A,B,C\) are \(16,2048,2048\), so all are nilpotent.
The definition permits \(B=C\).

The two points in each fiber have the same stabilizer,
\[
 B_w=\langle z_t:t\ne w\rangle\rtimes Q_w.
\]
Transitivity shows that these nine distinct groups are exactly
the conjugates of \(B\), accounting for base-group conjugation as well.
For each nonzero \(w\), \(H_w\) is an order-two reflection group
\(R_L\), where \(L=\langle w\rangle\).
The four lines give four distinct reflections: sign change of either
coordinate, interchange, and negative interchange.
Stabilizers of nonzero vectors on distinct lines intersect trivially.
It follows that all 81 ordered pairs of conjugates give precisely:
\[
\begin{array}{lll r}
\toprule
u,v & A\cap B_u\cap B_v & \text{order}&\text{pairs}\\
\midrule
u=v=o&H&8&1\\
\text{exactly one is }o,\ \text{other on }L&R_L&2&16\\
\text{both nonzero on }L&Z\times R_L&4&16\\
\text{nonzero on different lines}&Z&2&48\\
\bottomrule
\end{array}
\]
Membership in \(B_o\) kills \(Z\); all other \(B_w\)'s retain it.
The remaining condition is to fix the specified vectors in \(H\),
which proves every table entry.
There are ten distinct intersections, none trivial.
The five groups \(Z,R_L\) are exactly both the minimum-order and
the inclusion-minimal members. The reflections generate \(H\), so
\[
 \min_G(A,B,B)=\operatorname{Min}_G(A,B,B)=A.
\]

Finally \(F(G)=E\).
First \(O_2(Q)=1\): a normal \(2\)-subgroup centralizes the normal
translation subgroup \(\F_3^2\), whose centralizer in \(Q\) is itself.
Also \(C_G(E)=E\), because an affine part centralizing all independent
base flips fixes all nine indices, and the action is faithful.
For any normal nilpotent subgroup of \(G\), its characteristic odd
Sylow subgroups centralize \(E\), and hence are trivial.
It is therefore a \(2\)-group whose image in \(Q\) is trivial,
so lies in \(E\). Conversely \(E\) is normal abelian.
The reflection \(d\in A\) is outside \(E\), proving failure of all
assertions (a)--(c). In fact
\(A\cap B_o\cap B_{(0,1)}=\langle d\rangle\) is already a
minimum-order intersection with an outside witness.

The argument is elementary and determines the entire conjugacy and
intersection family; it is not an extrapolation from a sample.
The related affine stabilizer pattern has prior antecedents recorded
in \artifact{research/20.122-minimum-source-audit.md}.
The proof and independent enumeration packet are
\artifact{research/20.122-minimum-proof.md} and
\artifact{results/20.122-minimum-packet.json}.
